\textcolor{black}{\section{Procedimiento}}

Para determinar la carga eléctrica del electrón y su relación con la constante de Boltzmann se utilizo la ecuación Shockley para polarización directa como se muestras en la ecuación
\begin{equation}
	I_d = I_s \left[ \exp\left(\frac{eV_d}{nk_b T}\right) - 1 \right]
	\label{Shockley1}
\end{equation}

donde $I_s$ es la corriente de saturación inversa, $e$ es la carga del electrón $k_B$ el la constante de Bolztamann, T es la temperatura absoluta del diodo y $n$ es un factor entre 1 y 2, según el diodo 
opere en el modo difusivo (n≈ 1) o recombinativo (n≈2). Al desarrollar la  ecuación (\ref{Shockley1}) cuando $I_d >> I_s$ obtenemos:
\begin{equation}
    I_d \approx I_s \exp\left(\frac{eV_d}{nk_b T}\right)
    \label{shock2}
\end{equation}
Al aplicar la logaritmo a ambos lados se ajusta linealmente, la pendiente del ajuste nos muestra la relación $e/k_b$:
\begin{equation}
	\underbrace{\ln(I_d)}_{y} = \underbrace{\frac{e}{nk_b T}}_{m} \cdot \underbrace{V_d}_{x}  + \underbrace{ln(I_s)}_{b}
\end{equation}

Una vez obtenida la relación lineal, se procedió a obtener la dupla de datos ($V_d, I_d$) de manera experimental para esto se utilizo un modulo M05, que concite en un circuito eléctrico 